\chapter{Opis problemu i założeń}
\label{cha:problem}

\section{Opis problemu}

Problem polega na automatycznej identyfikacji użytkownika na podstawie danych biometrycznych w postaci \textbf{zdjęcia twarzy} oraz \textbf{nagrania głosu}. System musi:
\begin{enumerate}
    \item zarejestrować użytkowników (utworzyć wzorce biometryczne),
    \item porównać nowe dane próbne z bazą wszystkich użytkowników,
    \item wybrać najbardziej prawdopodobną tożsamość,
    \item podjąć decyzję o dostępie na podstawie progu akceptacji.
\end{enumerate}

W praktyce występują dwa rodzaje błędów~\cite{jain2004multimodal}:
\begin{itemize}
    \item \textbf{FAR} (\emph{ang. False Accept Rate}) --- fałszywa akceptacja, gdy system uznaje obcą osobę za zarejestrowanego użytkownika,
    \item \textbf{FRR} (\emph{ang. False Reject Rate}) --- fałszywe odrzucenie, gdy system odmawia dostępu prawdziwemu użytkownikowi.
\end{itemize}

Celem projektu jest minimalizacja obu metryk przy zachowaniu prostoty implementacji demonstracyjnej.

\section{Dane wejściowe}

System operuje na plikach przechowywanych lokalnie:
\begin{itemize}
    \item \textbf{Twarz:} obrazy \texttt{face\_*.jpg} lub \texttt{face\_*.png},
    \item \textbf{Głos:} nagrania \texttt{voice\_*.wav} lub \texttt{voice\_*.mp3}.
\end{itemize}

Dla każdego użytkownika zaleca się zebranie co najmniej trzech zdjęć i trzech nagrań (5--10~s mowy). Jako dane próbne wykorzystano pliki \texttt{face\_3} oraz \texttt{voice\_3}, różne od pozostałych wzorców rejestracyjnych, aby uniknąć zawyżonych wyników.

\section{Założenia działania}

Przyjęto następujące założenia:
\begin{itemize}
    \item jedna osoba na zdjęciu; twarz widoczna frontalnie lub lekko z boku,
    \item nagrania głosu zawierają wyłącznie mowę danej osoby, bez silnego szumu tła,
    \item identyfikacja odbywa się offline (brak wymogu chmury),
    \item baza użytkowników jest skończona i znana z góry (tryb 1:N),
    \item środowisko uruchomieniowe: Python~3.12, system Windows lub Linux.
\end{itemize}

\section{Ograniczenia}

Główne ograniczenia rozwiązania:
\begin{itemize}
    \item brak rejestracji i identyfikacji w czasie rzeczywistym z kamery lub mikrofonu,
    \item proste cechy głosu (MFCC), co daje mniejszą rozdzielczość niż modele głębokie,
    \item zależność modułu twarzy od biblioteki DeepFace i modelu Facenet,
    \item wrażliwość na jakość danych (oświetlenie, format plików audio),
    \item brak mechanizmów anty-spoofing (np.\ wykrywania zdjęcia zamiast żywej twarzy).
\end{itemize}
